\documentclass [11pt,a4paper,oneside,openany]{book} %Classe del documento, formato carta, singola facciata, apertura capitoli pag destra/sinista (ininfluente se impostato oneside). Formato libro
\usepackage[italian]{babel} %Lingua Documento (da impostare per la sillabazione)
\usepackage{graphicx} %Pacchetto necessario per la gestione delle immagini
\usepackage{xcolor,colortbl} %Colori
\usepackage{enumitem} % pacchetto per gestire le enumerazioni e resume
\usepackage{fancyhdr} %Pacchetto per la gestione accurata della pagina
\usepackage{setspace} %Pacchetto necessario per i comandi successivi onehalfspacing, singlespacing....
\usepackage{titlesec} % Pacchetto per le sezioni e sottosezioni
\usepackage{longtable} %Pacchetto per la gestione delle tabelle grandi
\usepackage[colorlinks=true]{hyperref} %Segnalibri nel pdf finale, sotto relativi parametri
\usepackage[T1]{fontenc}
\usepackage[utf8]{inputenc} %Pacchetti per le lettere accentate utf8
\usepackage{listings} %Pacchetto per formattare codice sorgente
\usepackage{mathtools} %Pacchetto matematica
\usepackage{float} %Pacchetto per il posizionamento delle figure
\usepackage{amsthm}

\hypersetup{
	bookmarksnumbered=true,
	linkcolor=black,
	citecolor=black,
	urlcolor=black,
}

\usepackage{geometry} % Dimensione pagina

\geometry{a4paper} % Formato carta
\addtolength{\textheight}{75pt} %Margini
\oddsidemargin 30pt

\setcounter{secnumdepth}{2}
\setcounter{tocdepth}{3}


\newcommand{\reqref}[1]{\textbf{R\ref{#1}}}
\newcommand{\goalref}[1]{\textbf{G\ref{#1}}}

\definecolor{darkgreen}{rgb}{0,0.6,0}
\lstset{frame=tb,
  language=Java,
  aboveskip=3mm,
  belowskip=3mm,
  showstringspaces=false,
  columns=flexible,
  basicstyle={\small\ttfamily},
  numbers=none,
  numberstyle=\tiny\color{gray},
  keywordstyle=\color{blue},
  commentstyle=\color{darkgreen},
  %stringstyle=\color{mauve},
  breaklines=true,
  breakatwhitespace=true,
  tabsize=3
}

\begin{document} %Inizio Documento

%%FRONTESPIZIO%%=========================================
\begin{titlepage}

 \begin{center} 
     \includegraphics[height=6cm]{logo.eps}\\
     \vspace{4em}
     {\Large \textsc{Software Engineering II}}\\
     \vspace{6em}
     {\LARGE \textbf{Travlendar+}}\\
     \vspace{3em}
     {\Large \textsc{Design}}\\
     \vspace{1em}
     {\Large \textsc{Document}}\\
 \end{center}
 
    \vskip 2cm
 
 	Authors:
 	\vspace{0.5em}
 	\begin{center}
      {\Large \textit{Edoardo D'Amico}}\\
      {\Large \textit{Gabbolini Giovanni}}\\
      {\Large \textit{Parroni Federico}}\\
    \end{center}

\vskip 2.0cm
\begin{center}
{\normalsize 1$^{st}$ November 2017}
\end{center}

\end{titlepage}

\newpage

%%INTESTAZIONI PAGINA%%==================================
\pagestyle{fancy}
\renewcommand{\chaptername}{Chapter}
\renewcommand{\chaptermark} [1]{\chaptername\ \thechapter.\ #1}{} 
\renewcommand{\chaptermark}[1]{\markboth{\thechapter.\ #1}{}} 
\renewcommand{\sectionmark}[1]{\markright{\thesection\ #1}}
\fancyhf{}
\fancyhead[LE,RO]{\bfseries\thepage} 
\fancyhead[LO,RE]{\bfseries\leftmark}
\fancypagestyle{plain}{%
\fancyhead{} % leva l'intestazione
\renewcommand{\headrulewidth}{0pt} % e la linea
}

%%INDICE%%==============================================
\renewcommand{\contentsname}{Indice}
\renewcommand{\listfigurename}{Elenco delle Figure}
\renewcommand{\listtablename}{Elenco delle Tabelle}
\tableofcontents

\linespread{1.5}\selectfont		%interlinea
%%CAPITOLI===============================================
%%CAPITOLO 1: Introduzione=======================================
\chapter{Introduction}

\section{Purpose}
The purpose of this document is to provide more technical details than the RASD about Travelendar+ system. This document is addressed to developers and aims to identify:
\begin{itemize}
\item The high level architecture
\item The design patterns
\item The main components and their interfaces provided one for another
\item The Runtime behavior
\end{itemize}

\section{Scope}

Travlendar+ is an application that aims to ease the management of the daily appointments of a registered user by providing a desired schedule for those. 
The application runs within the context of a restricted geographic area and thus should consider only the travel means available in this scope.
The application is able to interact with different data sources, in particular Travel Means APIs and User device GPS, and to arrange the appointments by exploiting these information.
Once the schedule is computed, the application gives to the user the possibility to buy tickets for the involved travel means, then he/she can run it and follow directions provided by the application. 

\section{Definitions, Acronyms, Synonims}

\subsection{Definitions}
\theoremstyle{definition}
\newtheorem{definition}{Definition}[section]
 
\begin{definition} \label{def:userdata}
The user data include:
\begin{itemize}
\item Schedules;
\item Appointments;
\item User parameters.
\end{itemize}

\end{definition}
  
\begin{definition} \label{def:device}
A device is a PC, a Tablet or a Smartphone in which run the last version of his O.S.;
\end{definition}

\begin{definition} \label{def:schedule}
A Schedule is a set of time-ordered and not overlapping appointments where their starting times are fixed and they're linked to each other by a path covered with a specific transportation mean;
\end{definition}

\begin{definition} \label{def:bestTravelMean}
The bestTravelMean is the travel mean that satisfies the constraints on the appointment and on the schedule and that optimize the chosen criteria.
\end{definition}


\subsection{Acronyms}
%Acronyms used in the text:
\begin{itemize}
\item API: Application Programming Interface;
\item CRUD: Create/Read/Update/Delete;
\end{itemize}

%\reqref{req:R1}
%\goalref{goal:G1}

\section{Revision history}

\section{Reference Documents}
RASD document

\section{Document structure}
\begin{itemize}
\item Introduction: this section introduces the design document. It contains a justification of his utility
\item Architecture Design: this section is divided into different parts:
\begin{enumerate}
\item Overview: this section explains the division in tiers of our application and lists further details that were not specified in the RASD.
\item Components View: this sections gives a global view of the components of the application, other than a brief description, and describes how they communicate
\item Deployment view: this section shows the components that must be deployed to have the application running correctly
\item Runtime view: sequence diagrams are represented in this section to show how the components interact and collaborate during the fullfilling of a specific task. Each task is associated to each goal specified on the RASD document
\item Component interfaces: the interfaces between the components are presented in this section
\item Selected architectural styles and patterns : this section explains the architectural and design choices taken during the creation of the application
\end{enumerate}
\item Algorithms Design: this section describes the most important algorithmic parts. Pseudo code is used in order to hide unnecessary implementation details and abstract from a specific programming language
\item Requirements Traceability: this section aims to explain how the decisions taken in the RASD are linked to design elements
\end{itemize}



 

\chapter{Architectural Design}

\section{Overview}
Travlendar+ has a multi-tier architecture. User can interact  directly with both version of the application: desktop and mobile. The schedule computation is performed locally on the device in which the application runs by querying the external APIs:

\begin{itemize}
\item Mapping Service API
\item Travel Means APIs
\item Wheather Forecasts APIs
\end{itemize}

For the last two ones, the communication happens through different interfaces (wrappers), in order to improve the capability of the system to be expanded. In this way we provide a common pattern for the data acquired from the APIs, so that new ones can be easily included without the need of drastic changes on the central core.

\section{Component View}

\section{Deployment View} 

\section{Runtime View}

\section{Component Interfaces}

\section{Selected architectural styles and patterns}

\section{Other Design Decisions}
\chapter{Algorithmic Design}
In this chapter are given the guidelines on how to implement the most important functionalities that the components of the system will offer. The pseudocode of the relevant method is shown. 

\section{Login}
The Login involves the application the application server and the database. The last one is located outside the application server. 
To login in to the application the following steps must be done:
\begin{enumerate}
\item the user presses the Login button in the Login view;
\item the Identity manager (Login method) sends a request to the Application Server for authentication
\item A query on the database is performed to check the validity of the user credential and a response to the client is returned;
\end{enumerate}

\begin{lstlisting}
//client side
Login(username, password)
	response = SendRequest("api/login", username, password)
	if response.isValid 
		token = response.GetToken
		Show(HomeView)
	else
		Show(ErrorMessage , "Invalid Credentials")

//server side
LoginRequest(username, password)
	result = database.query(username, password)
	if result == 1
		SendAuthenticationResponse (token)
	else
		sendAuthenticationError(error code)
\end{lstlisting}

\section{CreateSchedule}
\textbf{correggere formula}
The process of computing a schedule is composed by the following steps:

\begin{enumerate}
\item Create the predecessor matrix (P) in which every cell (i,j) contains:
\begin{equation}
    P_{ij}=
    \begin{cases}
      1, & if Appointment_i must precede Appointment_j \\
      0, & if Appointment_i must follows Appointment_j \\
      -1, & if Appointment_i can be scheduled both before or after appointment_j
    \end{cases}
  \end{equation}
\item Compute all the possible ordered arrangements of appointments with respect to the predecessor matrix
\item For each appointment in each arrangement, set the starting time according to travel time. This is estimated considering the euclidean distance and heuristics on the kind of travel mean between every pair of consecutive appointments in the arrangement, considering the bestTravelMean(\textbf{fare def});
\item Check the feasibility of the arrangements and discard the ones which have overlapping appointments;
\item Choose the most convenient one, according to the optimization criteria.
\item Call mapping service Api to fix the effective routes between the appointments
\end{enumerate}

\begin{lstlisting}
ComputeSchedule(wakeUpTime, startingLocation, appts, constraints, optCriteria)
	p=CalculatePredecessorMatrix(appts)
	a=CalculateArrangements(appts, p, 0, 1)
	SetStartingTime(a, startingTime, startingLocation, constraint, optCriteria)
	s=ChooseBestSchedule(a)
	MappingServiceRequest(s)
	return s

CalculatePredecessorMatrix(appts)
	p=new Matrix[appt.size,appt.size]
	for(i=0 .. appts.size)
		for(j=i+1 .. appts.size)
			a1=appts[i]
			a2=appts[j]
			if a1.deterministic && a2.deterministic
				if a1.startingTime < a2.startingTime
					pred[i,j]=1
				else
					pred[i,j]=0
			elseif a1.deterministic && !a2.deterministic
				if a1.startingTime < a2.timeSlot.start
					pred[i,j]=1
				elseif a1.endingTime > a2.timeSlot.end
					pred[i,j]=0
				else
					pred[i,j]=-1
			elseif !a1.deterministic && a2.deterministic
					if a1.timeslot.end < a2.endingTime
						pred[i,j]=1
					elseif a1.timeSlot.start > a2.startingTime
						pred[i,j]=0
					else
						pred[i,j]=-1
			elseif !a1.deterministic && !a2.deterministic
					if a1.timeSlot.end < a2.timeSlot.start
						pred[i,j]=1
					elseif a1.timeSlot.start > a2.timeSlot.end
						pred[i,j]=0
					else
						pred[i,j]=-1
	return p
		
		
CalculateArrangements(appts, p, curri, currj)
	arrangement=new List	
	for(i=curri .. appts.size-1)
		for(j=currj-1 .. appts.size)
			if p[i,j]==-1
				p0=p
				p0[i,j]=0
				CalculateArrangements(appts, p0, i, j)
				
				p1=p
				p1[i,j]=1
				CalculateArrangements(appts, p1, i, j)
				
				return
				
	a=ConvertPredMatrixToList(appts,p)
	arrangement.addLast(a)
	return arrangement

		 
ConvertPredMatrixToList(appts, p)
//converts a "-1 free" predecessor matrix to an ordered list of appointments


SetStartingTime(a, startingTime, startingLocation, constraint, optCriteria)
	for(arr in a)
		dummyStartingAppt = new appointment(startingTime, startingLocation, duration=0)
		arr.addFirst(dummyStartingAppt)
		for(i=1 .. arr.size)
			appt1=arr[i-1]
			appt2=arr[i]
			travelMean=getBestTravelMean(appt1, appt2, constraint, optCriteria)
			travelTime=travelMean.estimateTime(appt1, appt2)
			if appt2.deterministic 
				appt2.startingTravelTime = appt2.startingTime-travelTime
				appt2.travelMean=travelMean
				if appt1.endingTime > appt2.startingTravelTime
					error("schedule not feasible")
			else
				appt2.startingTravelTime = max(appt1.endingTime,appt2.timeSlot.start-travelTime)
				appt2.travelMean=travelMean
				if appt2.startingTravelTime > appt2.timeSlot.end
					error("schedule not feasible")
					
					
getBestTravelMean(appt1, appt2, constraint, optCriteria)
	l=getNotConstrainedTravelMeans(constraint)
	for(t in l)
		switch optCriteria
			case "MoneySpent"
				t.cost=estimateMoney(t, appt1, appt2)
			case "Time"
				t.cost=estimateTime(t, appt1, appt2)
			case "CarbonFootprint"
				t.cost=estimateCarbon(t, appt1, appt2)
	sortByCriteria(l, optCriteria)
	return l
	
ChooseBestSchedule(a)
	best=a[0]
	for(i = 0 .. a.size)
		sum=0
		for(appt in a[i])
			sum+=appt.travelMean.cost
		a[i].totalCost=sum
		if sum<a[0].totalCost
			best=a[i]
	return best
	
MappingServiceRequest(s)
	for (i=0 .. s.size-1)
		response=MappingServiceRequest(s[i], s[i+1])
		s[i].path=response.path
		s[i].startingTime=response.startingTime

\end{lstlisting}

\section{Registration}
The Registration involves the application the application server and the database. The last one is located outside the application server. 
To register in to the application the following steps must be done:
\begin{enumerate}
\item The user presses the Registration button in the Login view;
\item The Identity manager (Registration method) sends a request to the Application Server for authentication
\item A query on the database is performed to check the presence of the user credential and a confirmation email is sent
\item The user confirms the email by clicking on the designated link
\item The user's state on the database becomes confirmed
\end{enumerate}

\begin{lstlisting}
//client side
Register(username, password)
	response = SendRequest("api/registration", username, password)
	if !response.valid
		Show(ErrorMessage , "User Already Registered")
		

	
//server side
RegistrationRequest(username, password)
	result = database.query(username, password)
	if result == 0
		database.insertTuple(username, password)
		sendEmail(username)
	else
		sendAuthenticationError(error code)
		
EmailConfirmationRequest(username)
	result = database.modifyTuple(username, confirmed=true)
\end{lstlisting}


\section{Synchronize}
The synchronization in our system is the process that aims to keep the data consistent and updated between different devices. For example, when a user inserts an appointment on one of his devices, then the changes must be propagated to all his other devices, once the login is performed. 
The Synchronization involves the application, the application server and the database. The last one is located outside the application server. To sinchronize data across multiple devices two actions must be carried out:
\begin{itemize}
\item Upload of local data on the database when a single change on the  client's local data occours(Synchronize Upwards);
\item Download of data from the database, if an update is necessary, when login is performed(Synchronize Downwards).
\end{itemize}

\begin{lstlisting}
//client background processes

SynchronizeUpwards(changedData)
	update = false
	while !update
		response = SendRequest("api/sync/up", token, changedData)
		if response	 
			update = true
			
SynchronizeDownwards()
	newData=SendRequest("api/sync/down", token)
	localData=newData
		
		
//server side

SynchUpwardsRequest(token, changedData)
	userID = database.getUser(token)
	if changedData.action == insert
		//changedData is a newly inserted element
		database.insert(changedData, userID)
	else if changedData.action == edit
		database.update(changedData, userID)
	else
		database.delete(changedData, userID)
	sendResponse("syncresult", true)
	
SynchUpwardsRequest(token, changedData)
	userID = database.getUser(token)
	data=database.getUserData(userID)
	sendResponse(data)

\end{lstlisting}




		
\chapter{User Interface Design}
We refer to chapter 3.1.1 of the RASD, where we have given a complete overview of how the user interfaces will look like.
\chapter{Requirement Traceability}
In each section section we will describe how the requirements of each goal are mapped in the design elements described above.

\section{Goal 1}

\begin{itemize}

\item \textit{The system S.B.A to handle a registration phase in 
which the user will provide an e-mail and a password}

\begin{itemize}
\item Login Interface
\item Identity Manager
\item Authentication Provider
\item Database
\end{itemize}

\item \textit{The system S.B.A. to verify the e-mail given by the user by sending a confirmation e-mail to his address}

\begin{itemize}
\item Login Interface
\item Identity Manager
\item Authentication Provider
\item Database
\end{itemize}

\item \textit{The system should let the user to specify his parameters}

\begin{itemize}
\item User Account Interface
\end{itemize}

\end{itemize} 


\section{Goal 2}

\begin{itemize}

\item \textit{The system S.B.A to recognize a registered user given an e-mail and a password}

\begin{itemize}
\item Login Interface
\item Identity Manager
\item Authentication Provider
\item Database
\end{itemize}

\item \textit{The system S.B.A. to retrieve information from the user about his registration informations, i.e. his e-mail and password}

\begin{itemize}
\item Login Interface
\item Identity Manager
\item Authentication Provider
\item Database
\end{itemize}

\end{itemize}

\section{Goal 3}

\begin{itemize}

\item \textit{The system should be able to retrieve an e-mail address from the user}

\begin{itemize}
\item Login Interface
\end{itemize}

\item \textit{ The system should be able to send the password of a user to his e-mail given it}

\begin{itemize}
\item Authentication Provider
\item Database
\end{itemize}

\end{itemize}


\section{Goal 4}

\begin{itemize}

\item \textit{The system S.B.A. to retrieve information from the user about his appointments}

\begin{itemize}
\item Appointment CRUD interface.  \textbf{agg alle def}
\end{itemize}

\item \textit{The system S.B.A. to store an appointment in his memory}

\begin{itemize}
\item Appointment Manager
\end{itemize}

\end{itemize}


\section{Goal 5}

\begin{itemize}

\item \textit{The system should let the user change the parameters and the constraints of an inserted appointment}

\begin{itemize}
\item Appointment CRUD interface.  \textbf{agg alle def}
\item Appointment Manager
\end{itemize}

\item \textit{The system S.B.A to rewrite the appointment in his memory with his new parameters}

\begin{itemize}
\item Appointment Manager
\end{itemize}

\end{itemize}


\section{Goal 6}

\begin{itemize}

\item \textit{ Allow the user to set the constraints of the schedule}
\begin{itemize}
\item Schedule Interface
\end{itemize}

\item \textit{Allow the user to set the optimization criteria for the schedule}
\begin{itemize}
\item Schedule Interface
\end{itemize}

\item \textit{Allow the user to set the variables for the schedule}
\begin{itemize}
\item Schedule Interface
\end{itemize}

\item \textit{The system S.B.A. to gather information from external APIs about: 
	\begin{itemize}
		\item travel options with related travel option data; 
		\item weather forecast;
		\item strike days;
		\item delays.
	\end{itemize}}

\begin{itemize}
\item Mapping Service API Wrapper
\item Travel Mean API Wrapper
\item Weather Forecast API Wrapper
\end{itemize}

\item \textit{The system S.B.A. to select the best travel option according to the optimization criteria taking into account:
	\begin{itemize}
		\item user constraint;
		\item user parameters;
		\item travel option data; 
		\item weather forecast;
		\item information about strike day;
		\item information about delays.
	\end{itemize}}

\begin{itemize}
\item Scheduler
\end{itemize}

\item \textit{ The system S.B.A to store valid schedules requested by the user}
\begin{itemize}
\item Schedule Manager
\end{itemize}

\end{itemize}


\section{Goal 7}

\begin{itemize}

\item \textit{The system should let the user accept a schedule from the saved ones}
\begin{itemize}
\item Schedule List Interface
\item Schedule Manager
\end{itemize}

\end{itemize}

\section{Goal 8}

\begin{itemize}

\item \textit{The system S.B.A to book a travel mean through external API offered by third part application in which the user is signed}
\begin{itemize}
\item Schedule Manager
\item Travel Mean API Wrapper
\end{itemize}

\end{itemize}


\section{Goal 9}

\begin{itemize}

\item \textit{The system S.B.A to retrive the position of the user from his GPS}

\begin{itemize}
\item GPSAPI
\item Runtime Schedule Manager
\end{itemize}

\item \textit{ The system S.B.A. to retrive from an external API the directions to give to the user for reach the next appointment}

\begin{itemize}
\item Mapping Service API Wrapper
\item Runtime Schedule Manager
\end{itemize}

\item \textit{The system S.B.A. to retrive from an external API the Graphical representation of the path that will be travelled by the user}

\begin{itemize}
\item Mapping Service API Wrapper
\item Runtime Schedule Manager
\end{itemize}

\end{itemize}

\section{Goal 10}

\begin{itemize}

\item \textit{The system should be able to send notifications to the user}
\begin{itemize}
\item Schedule Manager
\end{itemize}

\item \textit{The system should be able to retrieve information about the availablity of shared travel means from external APIs without the user request}
\begin{itemize}
\item Mapping Service API Wrapper
\item Schedule Manager
\end{itemize}

\end{itemize}



\chapter{Implementation Integration and TestPlan}
In the next sections we will show the timeline of the implementation integration and testplan. This timeline is led by the functionalities, for each one we first implement the components needed, and then we proceed with the integration and with the test of the functionality.

\section{User Authentication}
Implemented and integrated Components:
\begin{itemize}
\item Authentication Provider
\item Database
\item Identity Manager
\end{itemize}
Functionality tested:
\begin{itemize}
\item User Registration
\item User Login
\item Recover credentials
\item SetUserParameters
\item GetUserParameters
\end{itemize}

\section{Appointments management}
Implemented and integrated components:
\begin{itemize}
\item Appointment manager
\end{itemize}
Functionality tested:
\begin{itemize}
\item Appointment creation
\item Appointment editing
\item Appointment deletion
\item Appointments retrieving
\end{itemize}

\section{External Api Data Retrieving}
Implemented and integrated components:
\begin{itemize}
\item Mapping Service API Wrapper
\item Travel Mean API Wrapper
\item Weather Forecast API Wrapper
\end{itemize}
Functionality tested:
\begin{itemize}
\item Weather forecast data retrieving
\item Strike day information retrieving
\item Travel option data retrieving
\item Get ticket price
\item Purchase ticket
\end{itemize}

\section{Static Schedule Management}
Implemented and integrated components:
\begin{itemize}
\item Scheduler
\item Schedule Manager
\end{itemize}
Functionality tested:
\begin{itemize}
\item Compute a schedule
\item Schedules retrieving
\item Schedule selection
\item Schedule editing
\item Schedule deletion
\item Schedule price updating
\end{itemize}

\section{Dynamic Schedule Management}
Implemented and integrated components:
\begin{itemize}
\item Runtime Schedule Manager
\end{itemize}
Functionality tested:
\begin{itemize}
\item Directions retrieving
\item Dynamic directions displaying
\end{itemize}

\section{Synchronization}
Implemented and integrated components:
\begin{itemize}
\item Synchronization provider
\item Synchronization Manager
\end{itemize}
Functionality tested:
\begin{itemize}
\item Synchronize
\end{itemize}




%%APPENDICI==============================================
\appendix
%%\input{appendiceA}

\end{document}